\documentclass[12pt]{article}
\usepackage[T1]{fontenc}
\usepackage[utf8]{inputenc}
\usepackage[brazil]{babel}
\usepackage[margin=1in]{geometry}
\setlength{\parindent}{0pt}
\begin{document}
\section*{Tema 1}
Processo: PEDILEF 2009.51.51.066212-3/ RJ\\
Situacao: Julgado\\
Ramo: DIREITO PREVIDENCIÁRIO\\
Questao: Saber qual a forma de cálculo da aposentadoria por invalidez, do auxílio-doença e pensões derivadas, em momento antecedente à edição da Lei n. 9.876/99.\\
Tese: O valor da aposentadoria por invalidez e do auxílio-doença, bem como das pensões destes derivados ou calculadas com base no art. 75, da Lei n. 8.213/91, será obtido, na forma do art. 29, II, do mesmo diploma, por meio da média aritmética simples dos 80\% maiores salários de contribuição, considerado todo o período contributivo, independentemente do momento de inscrição do segurado e do número de contribuições mensais do período contributivo.\\
Fonte: https://www.cjf.jus.br/phpdoc/virtus/
\end{document}
